\documentclass[../../main.tex]{subfiles}% 注意这里的文件路径不能用 ./main.tex ,否则用latexmk编译子文件会报错
\graphicspath{{\subfix{./image/}}{\subfix{../../image/}}} % 指定图片引用路径,后续可以直接使用图片文件名
% 如果图片存放在上级目录,则使用 ../../image/ 作为备用路径

% 例如:
% \begin{figure}[H]
% \centering
% \includegraphics[width=0.3\textwidth]{图.png}
% \caption{}
% \label{figure:图}
% \end{figure}
% 注意:上述\label{}一定要放在\caption{}之后,否则引用图片序号会只会显示??.

\begin{document}

\section{初值问题}

我们将在这一节中求解一维热方程初值问题
\begin{align}
\begin{cases}
\frac{\partial u}{\partial t} - a^2\frac{\partial^2 u}{\partial x^2} = f(x,t), & (x,t)\in\mathbb{R}\times\mathbb{R}_+, \\
u(x,0)=\varphi(x), & x\in\mathbb{R}.
\end{cases} \label{eq:heat-cauchy1d}
\end{align}
多维热方程初值问题的求解完全类似。在求解一维热方程初值问题 \eqref{eq:heat-cauchy1d} 之前我们先介绍 Fourier 变换和 Fourier 积分。

\subsection{Fourier 变换和 Fourier 积分}

Fourier 积分实际上是 Fourier 级数的一种极限形式。从 Fourier 级数，我们可以从形式上诱导出 Fourier 积分。

设 $f(x)\in C^1(\mathbb{R})$。对任意 $l>0$，$f(x)$ 在 $(-l,l)$ 可以展开为 Fourier 级数，且对 $x\in(-l,l)$，
\begin{align*}
f(x)=\frac{a_0}{2}+\sum_{k=1}^{\infty}\left(a_k\cos\frac{k\pi}{l}x+b_k\sin\frac{k\pi}{l}x\right),
\end{align*}
其中系数
\begin{align*}
a_k=\frac{1}{l}\int_{-l}^{l}f(\xi)\cos\frac{k\pi}{l}\xi\,\mathrm{d}\xi,\quad k=0,1,2,\cdots,\\
b_k=\frac{1}{l}\int_{-l}^{l}f(\xi)\sin\frac{k\pi}{l}\xi\,\mathrm{d}\xi,\quad k=1,2,\cdots.
\end{align*}
我们将系数 $a_k,b_k$ 的表达式代入上式，并利用三角函数和差化积公式得到
\begin{align*}
f(x)&=\frac{1}{2l}\int_{-l}^{l}f(\xi)\,\mathrm{d}\xi+\sum_{k=1}^{\infty}\frac{1}{l}\int_{-l}^{l}f(\xi)\cos\frac{k\pi}{l}(x-\xi)\,\mathrm{d}\xi \\
&=\frac{1}{2l}\int_{-l}^{l}f(\xi)\,\mathrm{d}\xi+\frac{1}{\pi}\sum_{k=1}^{\infty}\Delta\lambda_k\int_{-l}^{l}f(\xi)\cos\lambda_k(x-\xi)\,\mathrm{d}\xi,
\end{align*}
这里
\begin{align*}
\lambda_k=\frac{k\pi}{l},\qquad \Delta\lambda_k=\lambda_{k+1}-\lambda_k=\frac{\pi}{l}.
\end{align*}
设极限
\begin{align*}
\lim_{l\to+\infty}\int_{-l}^{l}|f(\xi)|\,\mathrm{d}\xi
\end{align*}
存在，则当 $l\to+\infty$ 时，
\begin{align*}
\frac{1}{2l}\int_{-l}^{l}f(\xi)\,\mathrm{d}\xi \longrightarrow 0,
\end{align*}
且形式上
\begin{align*}
\frac{1}{\pi}\sum_{k=1}^{\infty}\Delta\lambda_k\int_{-l}^{l}f(\xi)\cos\lambda_k(x-\xi)\,\mathrm{d}\xi
\longrightarrow \frac{1}{\pi}\int_0^{+\infty}\mathrm{d}\lambda\int_{-\infty}^{+\infty}f(\xi)\cos\lambda(x-\xi)\,\mathrm{d}\xi,
\end{align*}
因而我们可以猜测
\begin{align}
f(x)=\frac{1}{\pi}\int_0^{+\infty}\mathrm{d}\lambda\int_{-\infty}^{+\infty}f(\xi)\cos\lambda(x-\xi)\,\mathrm{d}\xi. \label{eq:fourier-informal}
\end{align}

\begin{definition}
设 $f(x)\in L^1(\mathbb{R})$，则无穷积分
\begin{align}
\hat{f}(\lambda)=\frac{1}{\sqrt{2\pi}}\int_{-\infty}^{+\infty}f(x)e^{-\mathrm{i}\lambda  x}\,\mathrm{d}x \label{eq:fourier-def}
\end{align}
有意义，称为 $f(x)$ 的 \textbf{Fourier 变换}，记为 $(f(x))^\wedge=\hat{f}(\lambda)$。
\end{definition}

\begin{lemma}[Riemann-Lebesgue 引理]
设 $f(x)\in C[a,b]$，则
\begin{align}
\lim_{n\to+\infty}\int_a^b f(x)\sin nx\,\mathrm{d}x=0. \label{eq:riemann-lebesgue}
\end{align}
\end{lemma}
\begin{remark}
另外，从复变函数我们知道如下有用的等式：
\begin{align}
\int_{-\infty}^{+\infty}\frac{\sin x}{x}\,\mathrm{d}x=\pi. \label{eq:sinc-integral}
\end{align}
我们将利用以上事实来证明 Fourier 积分定理。
\end{remark}

\begin{theorem}[Fourier 积分定理]
设 $f(x)\in L^1(\mathbb{R})\cap C^1(\mathbb{R})$，则
\begin{align}
f(x)=\lim_{N\to+\infty}\frac{1}{\sqrt{2\pi}}\int_{-N}^{N}\hat{f}(\lambda)e^{\mathrm{i}\lambda  x}\,\mathrm{d}\lambda. \label{eq:fourier-inversion}
\end{align}
公式\eqref{eq:fourier-inversion}称为\textbf{反演公式},右端的积分表示在 Cauchy 主值意义下的无穷积分，通常称为 \textbf{Fourier 逆变换}，记为 $(\hat{f}(\lambda))^\vee$。因此 Fourier 积分定理也可写成
\begin{align*}
(\hat{f})^\vee=f.
\end{align*}
也就是说，$L^1(\mathbb{R})\cap C^1(\mathbb{R})$ 中的函数作一次 Fourier 变换后，再作一次 Fourier 逆变换，又回到函数本身。
\end{theorem}
\begin{proof}
由于 $f(x)\in L^1(\mathbb{R})$，则对任意的 $\lambda\in\mathbb{R}$，无穷积分
\begin{align*}
\int_{-\infty}^{+\infty}f(x)e^{-\mathrm{i}\lambda  x}\,\mathrm{d}x
\end{align*}
一致收敛，且是 $\lambda$ 的连续函数。固定 $x\in\mathbb{R}$，我们利用上述无穷积分的一致收敛性，交换积分次序得
\begin{align*}
\frac{1}{\sqrt{2\pi}}\int_{-N}^{N}\hat{f}(\lambda)e^{\mathrm{i}\lambda  x}\,\mathrm{d}\lambda
&=\frac{1}{2\pi}\int_{-\infty}^{+\infty}f(\xi)\,\mathrm{d}\xi\int_{-N}^{N}e^{\mathrm{i}\lambda (x-\xi)}\,\mathrm{d}\lambda \\
&=\frac{1}{\pi}\int_{-\infty}^{+\infty}f(\xi)\frac{\sin N(x-\xi)}{x-\xi}\,\mathrm{d}\xi \\
&=\frac{1}{\pi}\int_{-\infty}^{+\infty}f(x+\eta)\frac{\sin N\eta}{\eta}\,\mathrm{d}\eta.
\end{align*}
设 $M$ 是一个待定正数。我们记
\begin{align*}
I_1=\frac{1}{\pi}\int_{-\infty}^{-M}f(x+\eta)\frac{\sin N\eta}{\eta}\,\mathrm{d}\eta,\\
I_2=\frac{1}{\pi}\int_{-M}^{M}f(x+\eta)\frac{\sin N\eta}{\eta}\,\mathrm{d}\eta,\\
I_3=\frac{1}{\pi}\int_{M}^{+\infty}f(x+\eta)\frac{\sin N\eta}{\eta}\,\mathrm{d}\eta,
\end{align*}
则
\begin{align*}
\frac{1}{\sqrt{2\pi}}\int_{-N}^{N}\hat{f}(\lambda)e^{\mathrm{i}\lambda  x}\,\mathrm{d}\lambda=I_1+I_2+I_3.
\end{align*}
由正弦函数的有界性我们估计 $I_1$ 和 $I_3$ 得到
\begin{align*}
|I_1|+|I_3|\leqslant \frac{1}{\pi M}\int_{-\infty}^{+\infty}|f(\eta)|\,\mathrm{d}\eta.
\end{align*}
接着我们估计 $I_2$：
\begin{align*}
I_2&=\frac{1}{\pi}\int_{-M}^{M}\frac{f(x+\eta)-f(x)}{\eta}\sin N\eta\,\mathrm{d}\eta+\frac{f(x)}{\pi}\int_{-M}^{M}\frac{\sin N\eta}{\eta}\,\mathrm{d}\eta \\
&=\frac{1}{\pi}\int_{-M}^{M}g(x,\eta)\sin N\eta\,\mathrm{d}\eta+\frac{f(x)}{\pi}\int_{-MN}^{MN}\frac{\sin\eta}{\eta}\,\mathrm{d}\eta,
\end{align*}
其中
\begin{align*}
g(x,\eta)=\int_0^1 f'(x+\tau\eta)\,\mathrm{d}\tau
\end{align*}
是关于 $\eta$ 的连续函数。对任意的 $\varepsilon>0$，先取定正数 $M$，使得
\begin{align*}
|I_1|+|I_3|<\frac{\varepsilon}{2},
\end{align*}
然后利用公式 \eqref{eq:riemann-lebesgue} 和 \eqref{eq:sinc-integral} 再取定 $N_0$，使得当 $N>N_0$ 时，
\begin{align*}
\left|\frac{1}{\pi}\int_{-M}^{M}g(x,\eta)\sin N\eta\,\mathrm{d}\eta\right|<\frac{\varepsilon}{4},
\end{align*}
和
\begin{align*}
\left|\frac{f(x)}{\pi}\int_{-MN}^{MN}\frac{\sin\eta}{\eta}\,\mathrm{d}\eta-f(x)\right|<\frac{\varepsilon}{4}.
\end{align*}
综上所述，我们得到，当 $N>N_0$ 时，
\begin{align*}
\left|\frac{1}{\sqrt{2\pi}}\int_{-N}^{N}\hat{f}(\lambda)e^{\mathrm{i}\lambda  x}\,\mathrm{d}\lambda-f(x)\right|\leqslant \varepsilon.
\end{align*}
从而完成定理的证明。
\end{proof}

\begin{proposition}[线性性质]
若 $f_i(x)\in L^1(\mathbb{R})$，$a_i\in\mathbb{C}$ $(i=1,2)$，则
\begin{align*}
(a_1f_1(x)+a_2f_2(x))^\wedge=a_1\hat{f}_1(\lambda)+a_2\hat{f}_2(\lambda).
\end{align*}
\end{proposition}

\begin{proposition}[微商性质]
若 $f(x),f'(x)\in L^1(\mathbb{R})\cap C(\mathbb{R})$，则
\begin{align*}
(f'(x))^\wedge=\mathrm{i}\lambda \hat{f}(\lambda).
\end{align*}
\end{proposition}

\begin{proof}
由于 $f'(x)\in C(\mathbb{R})\cap L^1(\mathbb{R})$，因此由 Newton-Leibniz 公式
\begin{align*}
f(x)=f(0)+\int_0^x f'(x)\,\mathrm{d}x
\end{align*}
我们知，极限
\begin{align*}
\lim_{x\to+\infty}f(x),\qquad \lim_{x\to-\infty}f(x)
\end{align*}
存在。又由于 $f(x)\in L^1(\mathbb{R})$，故
\begin{align*}
\lim_{|x|\to+\infty}f(x)=0.
\end{align*}
利用上式和分部积分公式得到
\begin{align*}
(f'(x))^\wedge=\frac{1}{\sqrt{2\pi}}\int_{-\infty}^{+\infty}f'(x)e^{-\mathrm{i}\lambda  x}\,\mathrm{d}x 
=\frac{1}{\sqrt{2\pi}}\mathrm{i}\lambda \int_{-\infty}^{+\infty}f(x)e^{-\mathrm{i}\lambda  x}\,\mathrm{d}x 
=\mathrm{i}\lambda \hat{f}(\lambda).
\end{align*}
\end{proof}

\begin{proposition}[乘多项式性质]
若 $f(x),xf(x)\in L^1(\mathbb{R})$，则
\begin{align*}
(xf(x))^\wedge=i\frac{\mathrm{d}}{\mathrm{d}\lambda}\hat{f}(\lambda).
\end{align*}
\end{proposition}

\begin{proof}
由于 $f(x),xf(x)\in L^1(\mathbb{R})$，故 $\hat{f}(\lambda)$ 是 $\lambda$ 的连续可微函数，且
\begin{align*}
\frac{\mathrm{d}}{\mathrm{d}\lambda}\hat{f}(\lambda)=\frac{1}{\sqrt{2\pi}}\int_{-\infty}^{+\infty}f(x)(-\mathrm{i}x)e^{-\mathrm{i}\lambda  x}\,\mathrm{d}x=-\mathrm{i}(xf(x))^\wedge.
\end{align*}
\end{proof}

\begin{corollary}
若 $f(x),f'(x),\cdots,f^{(m)}(x)\in L^1(\mathbb{R})\cap C(\mathbb{R})$，则
\begin{align*}
(f^{(m)}(x))^\wedge=(\mathrm{i}\lambda )^m\hat{f}(\lambda);
\end{align*}
若 $f(x),xf(x),\cdots,x^m f(x)\in L^1(\mathbb{R})$，则
\begin{align*}
(x^m f(x))^\wedge=\mathrm{i}^m\frac{\mathrm{d}^m}{\mathrm{d}\lambda^m}\hat{f}(\lambda).
\end{align*}
这里 $m$ 是大于或等于 $1$ 的正整数。
\end{corollary}

\begin{proposition}[平移性质]
对任意常数 $k$，若 $f(x)\in L^1(\mathbb{R})$，则
\begin{align*}
(f(x-k))^\wedge=e^{-\mathrm{i}\lambda  k}\hat{f}(\lambda).
\end{align*}
\end{proposition}

\begin{proof}
由定义得到
\begin{align*}
(f(x-k))^\wedge&=\frac{1}{\sqrt{2\pi}}\int_{-\infty}^{+\infty}f(x-k)e^{-\mathrm{i}\lambda  x}\,\mathrm{d}x \\
&=\frac{1}{\sqrt{2\pi}}\int_{-\infty}^{+\infty}f(y)e^{-\mathrm{i}\lambda (k+y)}\,\mathrm{d}y \\
&=e^{-\mathrm{i}\lambda  k}\hat{f}(\lambda).
\end{align*}
\end{proof}

\begin{proposition}[伸缩性质]
若 $f(x)\in L^1(\mathbb{R})$，则
\begin{align*}
(f(kx))^\wedge=\frac{1}{|k|}\hat{f}\left(\frac{\lambda}{k}\right),\quad k\neq 0.
\end{align*}
\end{proposition}

\begin{proof}
不妨设 $k>0$，由定义得到
\begin{align*}
(f(kx))^\wedge&=\frac{1}{\sqrt{2\pi}}\int_{-\infty}^{+\infty}f(kx)e^{-\mathrm{i}\lambda  x}\,\mathrm{d}x \\
&=\frac{1}{\sqrt{2\pi}}\int_{-\infty}^{+\infty}f(y)e^{-\mathrm{i}\lambda  \frac{y}{k}}\,\mathrm{d}\left(\frac{y}{k}\right) \\
&=\frac{1}{k}\hat{f}\left(\frac{\lambda}{k}\right).
\end{align*}
一般情形同理。
\end{proof}

\begin{proposition}[对称性质]
若 $f(x)\in L^1(\mathbb{R})$，则
\begin{align*}
(f(x))^\vee=\hat{f}(-\lambda).
\end{align*}
\end{proposition}

\begin{proposition}[卷积性质]
若 $f(x),g(x)\in L^1(\mathbb{R})$，那么 $f(x)$ 和 $g(x)$ 的卷积 $f*g$ 定义为
\begin{align*}
f*g(x)=\int_{-\infty}^{+\infty}f(x-t)g(t)\,\mathrm{d}t,
\end{align*}
则 $f*g(x)\in L^1(\mathbb{R})$，且
\begin{align*}
(f*g(x))^\wedge=\sqrt{2\pi}\hat{f}(\lambda)\hat{g}(\lambda).
\end{align*}
\end{proposition}

\begin{proof}
我们先证 $f*g(x)\in L^1(\mathbb{R})$。由 Fubini 定理得到
\begin{align*}
\int_{-\infty}^{+\infty}|f*g(x)|\,\mathrm{d}x
&\leqslant \int_{-\infty}^{+\infty}\mathrm{d}x\int_{-\infty}^{+\infty}|f(x-t)g(t)|\,\mathrm{d}t \\
&=\int_{-\infty}^{+\infty}|g(t)|\,\mathrm{d}t\int_{-\infty}^{+\infty}|f(x-t)|\,\mathrm{d}x \\
&=\int_{-\infty}^{+\infty}|g(t)|\,\mathrm{d}t\cdot\int_{-\infty}^{+\infty}|f(x)|\,\mathrm{d}x,
\end{align*}
从而 $f*g\in L^1(\mathbb{R})$。

又由 Fubini 定理，我们知
\begin{align*}
(f*g(x))^\wedge&=\frac{1}{\sqrt{2\pi}}\int_{-\infty}^{+\infty}e^{-\mathrm{i}\lambda  x}\mathrm{d}x\int_{-\infty}^{+\infty}f(x-t)g(t)\,\mathrm{d}t \\
&=\frac{1}{\sqrt{2\pi}}\int_{-\infty}^{+\infty}g(t)e^{-\mathrm{i}\lambda  t}\left[\int_{-\infty}^{+\infty}f(x-t)e^{-\mathrm{i}\lambda (x-t)}\,\mathrm{d}x\right]\mathrm{d}t \\
&=\sqrt{2\pi}\hat{f}(\lambda)\hat{g}(\lambda).
\end{align*}
\end{proof}

\begin{example}
设
\begin{align*}
f_1(x)=\begin{cases}
1, & |x|\leqslant A,\\
0, & |x|>A
\end{cases}\quad (A>0).
\end{align*}
求 $\hat{f}_1(\lambda)$。
\end{example}

\begin{solution}
由定义，有
\begin{align*}
\hat{f}_1(\lambda)=\frac{1}{\sqrt{2\pi}}\int_{-A}^{A}e^{-\mathrm{i}\lambda  x}\,\mathrm{d}x=\sqrt{\frac{2}{\pi}}\frac{\sin\lambda A}{\lambda}.
\end{align*}
\end{solution}

此例可以用于一维波动方程初值问题的求解。

\begin{example}
设
\begin{align*}
f_2(x)=\begin{cases}
e^{-x}, & x\geqslant 0,\\
0, & x<0.
\end{cases}
\end{align*}
求 $\hat{f}_2(\lambda)$。
\end{example}

\begin{solution}
由定义，有
\begin{align*}
\hat{f}_2(\lambda)=\frac{1}{\sqrt{2\pi}}\int_0^{+\infty}e^{-(1+\mathrm{i}\lambda )x}\,\mathrm{d}x=\frac{1}{\sqrt{2\pi}(1+\mathrm{i}\lambda )}.
\end{align*}
\end{solution}

\begin{example}
设 $f_3(x)=e^{-|x|}$。求 $\hat{f}_3(\lambda)$。
\end{example}

\begin{solution}
由于 $f_3(x)=f_2(x)+f_2(-x)$，利用线性性质和伸缩性质则有
\begin{align*}
\hat{f}_3(\lambda)&=\hat{f}_2(\lambda)+(f_2(-x))^\wedge 
=\hat{f}_2(\lambda)+\hat{f}_2(-\lambda) \\
&=\frac{1}{\sqrt{2\pi}(1+\mathrm{i}\lambda )}+\frac{1}{\sqrt{2\pi}(1-\mathrm{i}\lambda )} 
=\frac{2}{\sqrt{2\pi}(1+\lambda^2)}.
\end{align*}
\end{solution}

\begin{example}
设 $f_4(x)=e^{-x^2}$。求 $\hat{f}_4(\lambda)$。
\end{example}

\begin{solution}
由分部积分公式和乘多项式性质我们得到，当 $\lambda\neq 0$ 时，
\begin{align*}
\hat{f}_4(\lambda)&=\frac{1}{\sqrt{2\pi}}\int_{-\infty}^{+\infty}e^{-x^2}e^{-\mathrm{i}\lambda  x}\,\mathrm{d}x \\
&=-\frac{1}{\sqrt{2\pi}}\cdot\frac{1}{\mathrm{i}\lambda }\left(e^{-x^2}e^{-\mathrm{i}\lambda  x}\Big|_{-\infty}^{+\infty}+2\int_{-\infty}^{+\infty}xe^{-x^2}e^{-\mathrm{i}\lambda  x}\,\mathrm{d}x\right) \\
&=\frac{2i}{\lambda}(xf_4(x))^\wedge=-\frac{2}{\lambda}\cdot\frac{\mathrm{d}}{\mathrm{d}\lambda}\hat{f}_4(\lambda).
\end{align*}
注意到恒等式
\begin{align*}
\int_{-\infty}^{+\infty}e^{-x^2}\,\mathrm{d}x=\sqrt{\pi},
\end{align*}
我们导出 $\hat{f}_4(\lambda)$ 作为 $\lambda$ 的函数满足下列常微分方程初值问题
\begin{align*}
\begin{cases}
\frac{\mathrm{d}}{\mathrm{d}\lambda}\hat{f}_4(\lambda)=-\frac{\lambda}{2}\hat{f}_4(\lambda),\\
\hat{f}_4(0)=\frac{1}{\sqrt{2}}.
\end{cases}
\end{align*}
求解上述常微分方程初值问题得 $\hat{f}_4(\lambda)=\frac{1}{\sqrt{2}}e^{-\lambda^2/4}$。
\end{solution}

\begin{example}
设
\begin{align*}
f_5(x)=e^{-Ax^2}\quad (A>0).
\end{align*}
求 $\hat{f}_5(\lambda)$。
\end{example}

\begin{solution}
利用伸缩性质得到
\begin{align*}
\hat{f}_5(\lambda)=(f_4(\sqrt{A}x))^\wedge=\frac{1}{\sqrt{A}}\hat{f}_4\left(\frac{\lambda}{\sqrt{A}}\right)
=\frac{1}{\sqrt{2A}}e^{-\frac{\lambda^2}{4A}}.
\end{align*}
\end{solution}

\begin{definition}
设 $f(x)=f(x_1,x_2,\dots,x_n)\in L^1(\mathbb{R}^n)$，则无穷积分
\begin{align}
\hat{f}(\lambda)=\frac{1}{(\sqrt{2\pi})^n}\int_{\mathbb{R}^n}f(x)e^{-\mathrm{i}\lambda \cdot x}\,\mathrm{d}x \label{eq:fourier-multidim}
\end{align}
有意义，称为 $f(x)$ 的 Fourier 变换，这里 $\lambda=(\lambda_1,\lambda_2,\dots,\lambda_n)\in\mathbb{R}^n$，记为 $(f(x))^\wedge=\hat{f}(\lambda)$。
\end{definition}

\begin{theorem}[反演公式]
设 $f(x)\in L^1(\mathbb{R}^n)\cap C^1(\mathbb{R}^n)$，则
\begin{align}
(\hat{f}(\lambda))^\vee=\lim_{N\to+\infty}\frac{1}{(\sqrt{2\pi})^n}\int_{|\lambda|\leqslant N}\hat{f}(\lambda)e^{\mathrm{i}\lambda \cdot x}\,\mathrm{d}\lambda=f(x), \label{eq:fourier-inv-multidim}
\end{align}
其中 $x=(x_1,x_2,\dots,x_n)\in\mathbb{R}^n$。$(\hat{f}(\lambda))^\vee$ 表示在 Cauchy 主值意义下的无穷积分，称为 $\hat{f}(\lambda)$ 的 Fourier 逆变换。换句话说，
\begin{align*}
(\hat{f})^\vee=f.
\end{align*}
\end{theorem}

对于多维 Fourier 变换，容易证明有类似于一维 Fourier 变换的性质成立。

\begin{proposition}[线性性质]
若 $f_j(x)\in L^1(\mathbb{R}^n)$，$a_j\in\mathbb{C}$ $(j=1,2)$，则
\begin{align*}
(a_1f_1(x)+a_2f_2(x))^\wedge=a_1\hat{f}_1(\lambda)+a_2\hat{f}_2(\lambda).
\end{align*}
\end{proposition}

\begin{proposition}[微商性质]
若 $f(x),D f(x)\in L^1(\mathbb{R}^n)\cap C(\mathbb{R}^n)$，则
\begin{align*}
\left(\frac{\partial f}{\partial x_j}\right)^\wedge=\mathrm{i}\lambda _j\hat{f}(\lambda),\quad j=1,2,\dots,n.
\end{align*}
\end{proposition}

\begin{proposition}[乘多项式性质]
若 $f(x),x_j f(x)\in L^1(\mathbb{R}^n)$，$j=1,2,\dots,n$，则
\begin{align*}
(x_j f(x))^\wedge=i\frac{\partial}{\partial\lambda_j}\hat{f}(\lambda).
\end{align*}
\end{proposition}

\begin{remark}
作为上述两个性质的推论，若 $f(x),D f(x),\cdots,D^k f(x)\in L^1(\mathbb{R}^n)\cap C(\mathbb{R}^n)$，则
\begin{align*}
(D^\alpha f(x))^\wedge=(\mathrm{i}\lambda )^\alpha\hat{f}(\lambda);
\end{align*}
若 $x^\alpha f(x)\in L^1(\mathbb{R}^n)$，则
\begin{align*}
(x^\alpha f(x))^\wedge=i^{|\alpha|}D^\alpha\hat{f}(\lambda).
\end{align*}
这里 $\alpha=(\alpha_1,\alpha_2,\dots,\alpha_n)$ 是阶数小于或等于 $k$ 的多重指标，$x^\alpha=x_1^{\alpha_1}x_2^{\alpha_2}\cdots x_n^{\alpha_n}$。
\end{remark}

\begin{proposition}[平移性质]
对任意固定的 $x_0\in\mathbb{R}^n$，若 $f(x)\in L^1(\mathbb{R}^n)$，则
\begin{align*}
(f(x-x_0))^\wedge=e^{-\mathrm{i}\lambda \cdot x_0}\hat{f}(\lambda).
\end{align*}
\end{proposition}

\begin{proposition}[伸缩性质]
若 $f(x)\in L^1(\mathbb{R}^n)$，则
\begin{align*}
(f(kx))^\wedge=\frac{1}{|k|}\hat{f}\left(\frac{\lambda}{k}\right),\quad k\neq 0.
\end{align*}
\end{proposition}

\begin{proposition}[对称性质]
若 $f(x)\in L^1(\mathbb{R}^n)$，则
\begin{align*}
(f(x))^\vee=\hat{f}(-\lambda).
\end{align*}
\end{proposition}

\begin{proposition}[卷积性质]
若 $f(x),g(x)\in L^1(\mathbb{R}^n)$，那么 $f(x)$ 和 $g(x)$ 的卷积 $f*g$ 定义为
\begin{align*}
f*g(x)=\int_{\mathbb{R}^n}f(x-y)g(y)\,\mathrm{d}y,
\end{align*}
则 $f*g(x)\in L^1(\mathbb{R}^n)$ 且
\begin{align*}
(f*g(x))^\wedge=(\sqrt{2\pi})^n\hat{f}(\lambda)\hat{g}(\lambda).
\end{align*}
\end{proposition}

同时我们还有如下结论。

\begin{proposition}[分离变量性质]
若 $f(x)=f_1(x_1)f_2(x_2)\cdots f_n(x_n)$，其中 $f_j(x_j)\in L^1(\mathbb{R})$，则
\begin{align*}
\hat{f}(\lambda)=\prod_{i=1}^{n}\hat{f}_i(\lambda_i).
\end{align*}
这里 $x=(x_1,x_2,\cdots,x_n)\in\mathbb{R}^n$，$\lambda=(\lambda_1,\lambda_2,\cdots,\lambda_n)\in\mathbb{R}^n$。
\end{proposition}

\begin{proof}
由定义将重积分化为累次积分即得证。
\end{proof}

\begin{example}
设 $f_6(x)=e^{-A|x|^2}$ ($A>0$)。求 $\hat{f}_6(\lambda)$。
\end{example}

\begin{solution}
事实上，由分离变量性质我们得到
\begin{align*}
\hat{f}_6(\lambda)=\prod_{i=1}^{n}(e^{-Ax_i^2})^\wedge=\prod_{i=1}^{n}\frac{1}{\sqrt{2A}}e^{-\frac{\lambda_i^2}{4A}}=\frac{1}{(2A)^{\frac{n}{2}}}e^{-\frac{|\lambda|^2}{4A}}.
\end{align*}
\end{solution}

\subsection{初值问题和基本解}

现在我们利用 Fourier 变换来求解初值问题 \eqref{eq:heat-cauchy1d}。对方程和初始条件两边关于 $x$ 作 Fourier 变换，利用线性性质和微商性质得到
\begin{align*}
\begin{cases}
\frac{\mathrm{d}\hat{u}}{\mathrm{d}t}+a^2\lambda^2\hat{u}=\hat{f}(\lambda,t),\\
\hat{u}(\lambda,0)=\hat{\varphi}(\lambda),
\end{cases}
\end{align*}
其中 $\hat{u}=\hat{u}(\lambda,t)$ 是解 $u(x,t)$ 关于 $x$ 的 Fourier 变换。解这个常微分方程的初值问题得
\begin{align*}
\hat{u}(\lambda,t)=\hat{\varphi}e^{-a^2\lambda^2 t}+\int_0^t \hat{f}(\lambda,\tau)e^{-a^2\lambda^2(t-\tau)}\,\mathrm{d}\tau.
\end{align*}
然后对上式两边求 Fourier 逆变换得
\begin{align*}
u(x,t)=(\hat{u}(\lambda,t))^\vee 
=(\hat{\varphi}e^{-a^2\lambda^2 t})^\vee+\int_0^t(\hat{f}(\lambda,\tau)e^{-a^2\lambda^2(t-\tau)})^\vee\,\mathrm{d}\tau.
\end{align*}
利用伸缩性质的结论，我们得到
\begin{align*}
e^{-a^2\lambda^2 t}=(g(x,t))^\wedge,
\end{align*}
其中
\begin{align*}
g(x,t)=\frac{1}{a\sqrt{2t}}e^{-\frac{x^2}{4a^2t}},
\end{align*}
从而
\begin{align*}
(\hat{\varphi}e^{-a^2\lambda^2 t})^\vee=(\hat{\varphi}\hat{g})^\vee=\frac{1}{\sqrt{2\pi}}\varphi*g
=\frac{1}{2a\sqrt{\pi t}}\int_{-\infty}^{+\infty}\varphi(\xi)e^{-\frac{(x-\xi)^2}{4a^2t}}\,\mathrm{d}\xi.
\end{align*}
同理
\begin{align*}
(\hat{f}(\lambda,\tau)e^{-a^2\lambda^2(t-\tau)})^\vee
=\frac{1}{2a\sqrt{\pi(t-\tau)}}\int_{-\infty}^{+\infty}f(\xi,\tau)e^{-\frac{(x-\xi)^2}{4a^2(t-\tau)}}\,\mathrm{d}\xi.
\end{align*}
于是我们得到
\begin{align}
u(x,t)=\int_{-\infty}^{+\infty}K(x-\xi,t)\varphi(\xi)\,\mathrm{d}\xi+\int_0^t\mathrm{d}\tau\int_{-\infty}^{+\infty}K(x-\xi,t-\tau)f(\xi,\tau)\,\mathrm{d}\xi, \label{eq:poisson-1d}
\end{align}
这里
\begin{align}
K(x,t)=\begin{cases}
\frac{1}{2a\sqrt{\pi t}}e^{-\frac{x^2}{4a^2t}}, & t>0,\\
0, & t\leqslant 0
\end{cases} \label{eq:poisson-kernel1d}
\end{align}
称为\textbf{Poisson 核}。

表达式 \eqref{eq:poisson-1d} 称为 \textbf{Poisson 公式}，$\Gamma(x,t;\xi,\tau)=K(x-\xi,t-\tau)$ 称为\textbf{热方程的基本解}。热方程的基本解也可以如下定义。

\begin{definition}
对于任意 $(\xi,\tau)\in\mathbb{R}_+^2=\mathbb{R}\times\mathbb{R}_+$，若函数 $u(x,t)\in L_{\mathrm{loc}}^1(\mathbb{R}_+^2)\cap C(\mathbb{R}_+^2\setminus(\xi,\tau))$ 且在广义函数意义下满足方程和初始条件
\begin{align*}
\begin{cases}
\frac{\partial u}{\partial t}-a^2\frac{\partial^2 u}{\partial x^2}=\delta(x-\xi,t-\tau), & (x,t)\in\mathbb{R}\times\mathbb{R}_+,\\
u(x,0)=0, & x\in\mathbb{R},
\end{cases}
\end{align*}
则称之为\textbf{热方程的基本解}，且记为 $\Gamma(x,t;\xi,\tau)$。这里 $\delta(x-\xi,t-\tau)$ 是 Dirac 测度。
\end{definition}

基本解的物理意义如下：考虑两端均在无穷远点，侧表面绝热的均匀细杆。在 $t=\tau$ 时刻，在 $x=\xi$ 处放置一个瞬时单位点热源，那么由这个瞬时单位点热源在杆上产生的温度分布实际上就是基本解 $\Gamma(x,t;\xi,\tau)$。因此基本解 $\Gamma(x,t;\xi,\tau)$ 也称为\textbf{点源函数}。

\begin{proposition}
当 $t>\tau$ 时，$\Gamma(x,t;\xi,\tau)>0$。
\end{proposition}

\begin{proposition}
$\Gamma(x,t;\xi,\tau)=\Gamma(\xi,t;x,\tau)$。
\end{proposition}

\begin{proposition}
当 $t>\tau$，$x\in\mathbb{R}$ 时，
\begin{align*}
\int_{-\infty}^{+\infty}\Gamma(x,t;\xi,\tau)\,\mathrm{d}\xi=1.
\end{align*}
\end{proposition}

\begin{proof}
当 $t>\tau$ 时，作变量替换 $\eta=(\xi-x)/(2a\sqrt{t-\tau})$，则
\begin{align*}
\int_{-\infty}^{+\infty}\Gamma(x,t;\xi,\tau)\,\mathrm{d}\xi
=\int_{-\infty}^{+\infty}\frac{1}{2a\sqrt{\pi(t-\tau)}}e^{-\frac{(x-\xi)^2}{4a^2(t-\tau)}}\,\mathrm{d}\xi 
=\frac{1}{\sqrt{\pi}}\int_{-\infty}^{+\infty}e^{-\eta^2}\,\mathrm{d}\eta=1.
\end{align*}
\end{proof}

\begin{proposition}
当 $t>\tau$，$x\in\mathbb{R}$ 时，$\Gamma(x,t;\xi,\tau)$ 关于所有自变量无穷次连续可微，且满足方程
\begin{align*}
\frac{\partial\Gamma}{\partial t}-a^2\frac{\partial^2\Gamma}{\partial x^2}=0,\qquad \frac{\partial\Gamma}{\partial\tau}+a^2\frac{\partial^2\Gamma}{\partial\xi^2}=0.
\end{align*}
\end{proposition}

\begin{proposition}
当 $t>\tau$ 时，成立估计式
\begin{align*}
|\Gamma(x,t;\xi,\tau)|\leqslant \frac{1}{2a\sqrt{\pi(t-\tau)}}.
\end{align*}
\end{proposition}

\begin{proposition}
若 $\varphi\in C(\mathbb{R})\cap L^\infty(\mathbb{R})$，即 $\varphi$ 是 $\mathbb{R}$ 上连续的有界函数，则对于任意 $x\in\mathbb{R}$，成立
\begin{align*}
\lim_{t\to 0^+}\int_{-\infty}^{+\infty}\Gamma(x,t;\xi,0)\varphi(\xi)\,\mathrm{d}\xi=\varphi(x).
\end{align*}
\end{proposition}

\begin{proof}
作变量替换 $\eta=(\xi-x)/(2a\sqrt{t})$，则
\begin{align*}
\int_{-\infty}^{+\infty}\Gamma(x,t;\xi,0)\varphi(\xi)\,\mathrm{d}\xi
=\frac{1}{\sqrt{\pi}}\int_{-\infty}^{+\infty}e^{-\eta^2}\varphi(x+2a\sqrt{t}\eta)\,\mathrm{d}\eta.
\end{align*}
利用 $\varphi$ 的有界性和连续性，上述积分在 $t\in\mathbb{R}_+$ 上关于 $t$ 一致收敛。因此
\begin{align*}
\lim_{t\to 0^+}\int_{-\infty}^{+\infty}\Gamma(x,t;\xi,0)\varphi(\xi)\,\mathrm{d}\xi
&=\frac{1}{\sqrt{\pi}}\int_{-\infty}^{+\infty}e^{-\eta^2}\lim_{t\to 0^+}\varphi(x+2a\sqrt{t}\eta)\,\mathrm{d}\eta \\
&=\varphi(x)\frac{1}{\sqrt{\pi}}\int_{-\infty}^{+\infty}e^{-\eta^2}\,\mathrm{d}\eta=\varphi(x).
\end{align*}
\end{proof}

\begin{theorem}
设 $\varphi(x)\in C(\mathbb{R})\cap L^\infty(\mathbb{R})$，$f(x,t)\equiv 0$，则 Poisson 公式 \eqref{eq:poisson-1d} 表示的函数 $u(x,t)$ 是初值问题 \eqref{eq:heat-cauchy1d} 的有界 $C^{2,1}(\mathbb{R}_+^2)$ 解，其中函数 $u(x,t)$ 在极限的意义下取初值 $\varphi(x)$。
\end{theorem}
\begin{remark}
定理中关于 $\varphi(x)$ 有界的假设可以改进为
\begin{align*}
|\varphi(x)|\leqslant Me^{Ax^2},\quad x\in\mathbb{R},
\end{align*}
这里 $A$ 和 $M$ 是正常数。此时 Poisson 公式 \eqref{eq:poisson-1d} 表示的函数 $u(x,t)$ 在区域 $\mathbb{R}\times(0,(4a^2A)^{-1})$ 上仍然是初值问题 \eqref{eq:heat-cauchy1d} 的 $C^\infty$ 解。
\end{remark}
\begin{proof}
当 $(x,t)\in\mathbb{R}_+^2$ 时，
\begin{align*}
|u(x,t)|&\leqslant \frac{1}{2a\sqrt{\pi t}}\int_{-\infty}^{+\infty}|\varphi(\xi)|e^{-\frac{(x-\xi)^2}{4a^2t}}\,\mathrm{d}\xi \\
&\leqslant \sup_{\mathbb{R}}|\varphi(\xi)|\frac{1}{\sqrt{\pi}}\int_{-\infty}^{+\infty}e^{-\eta^2}\,\mathrm{d}\eta \\
&\leqslant \sup_{\mathbb{R}}|\varphi(\xi)|.
\end{align*}
另外，固定任何一点 $(x_0,t_0)\in\mathbb{R}_+^2$，我们在其邻域 $(x_0-t_0,x_0+t_0)\times(\frac{t_0}{2},\frac{3t_0}{2})$ 内考虑 $u(x,t)$ 的微分性质。因此，$\frac{1}{\sqrt{t}}$ 在此邻域上无穷次可微，而 $e^{-\frac{(x-\xi)^2}{4a^2t}}$ 与任何次数的多项式 $P(\xi)$ 的乘积在 $\mathbb{R}$ 上可积，根据微积分的一致收敛定理我们知道 Poisson 公式 \eqref{eq:poisson-1d} 中的 $u(x,t)$ 关于 $x,t$ 求任意阶偏导数可以与关于 $\xi$ 求积分交换次序。由此我们得到 $u(x,t)$ 是 $C^\infty(\mathbb{R}_+^2)$ 中的函数，而且
\begin{align*}
\frac{\partial u(x,t)}{\partial t}-a^2\frac{\partial^2 u(x,t)}{\partial x^2}
=\int_{-\infty}^{+\infty}\left[\frac{\partial K(x-\xi,t)}{\partial t}-a^2\frac{\partial^2 K(x-\xi,t)}{\partial x^2}\right]\varphi(\xi)\,\mathrm{d}\xi=0.
\end{align*}
从而证明 $u(x,t)$ 满足热方程。

接着我们证明 $u(x,t)$ 满足初值条件。在 $u(x,t)$ 的积分表达式中作变量替换 $\eta=(\xi-x)/(2a\sqrt{t})$，则
\begin{align*}
u(x,t)=\frac{1}{\sqrt{\pi}}\int_{-\infty}^{+\infty}e^{-\eta^2}\varphi(x+2a\sqrt{t}\eta)\,\mathrm{d}\eta.
\end{align*}
利用 $\varphi$ 的有界性和连续性，上述积分在 $(x,t)\in\mathbb{R}\times\mathbb{R}_+$ 上关于 $x,t$ 一致收敛。因此
\begin{align*}
\lim_{\substack{x\to x_0\\ t\to 0^+}}u(x,t)
&=\frac{1}{\sqrt{\pi}}\int_{-\infty}^{+\infty}e^{-\eta^2}\varphi(x_0)\,\mathrm{d}\eta=\varphi(x_0).
\end{align*}
至此完成定理的证明。
\end{proof}

\begin{proposition}[奇偶性和周期性]
设 $\varphi$ 是奇（或偶，或周期为 $T$ 的）函数，则解 $u(x,t)$ 也是 $x$ 的奇（或偶，或周期为 $T$ 的）函数。
\end{proposition}

\begin{proposition}[无限传播速度]
设无穷长杆的初始温度 $\varphi(x)$ 在小段 $(x_0-\delta,x_0+\delta)$ 大于零，且在其他地方 $\mathbb{R}\setminus(x_0-\delta,x_0+\delta)$ 等于零，则杆在 $t>0$ 时刻，在任何一点 $x$ 的温度
\begin{align*}
u(x,t)=\int_{-\infty}^{+\infty}K(x-\xi,t)\varphi(\xi)\,\mathrm{d}\xi>0.
\end{align*}
\end{proposition}
\begin{remark}
此性质说明，顷刻之间，杆的热量传递到杆上的任何一点，且离 $x_0$ 近的点，温度会升高多一些；而离 $x_0$ 远的点，温度会升高少一些。当然，这个结论与线性的 Fourier 热传导定律有关。
\end{remark}

\begin{proposition}[无穷次可微性]
若 $\varphi(x)\in L^\infty(\mathbb{R})\cap C(\mathbb{R})$，不管 $\varphi(x)$ 是否可微，当 $t>0$ 时，解 $u(x,t)$ 在 $\mathbb{R}\times\mathbb{R}_+$ 上必无穷次可微。
\end{proposition}

事实上，当 $t>0$ 时，对于任意点 $x\in\mathbb{R}$，
\begin{align*}
\frac{\partial^{k+l}u}{\partial x^k\partial t^l}
=\int_{-\infty}^{+\infty}\frac{\partial^{k+l}}{\partial x^k\partial t^l}K(x-\xi,t)\varphi(\xi)\,\mathrm{d}\xi.
\end{align*}
这是因为对于任何正数 $\delta$，上述积分在区域 $\mathbb{R}\times(\delta,+\infty)$ 一致收敛，因而可以将求微商与求积分交换次序。

同样我们可以利用多维 Fourier 变换来求解多维热方程初值问题
\begin{align}
\begin{cases}
\frac{\partial u}{\partial t}-a^2\Delta u=f(x,t), & (x,t)\in\mathbb{R}_+^{n+1}=\mathbb{R}^n\times\mathbb{R}_+,\\
u(x,0)=\varphi(x), & x\in\mathbb{R}^n.
\end{cases} \label{eq:heat-cauchy-nd}
\end{align}
我们可以得到 Poisson 公式
\begin{align}
u(x,t)=\int_{\mathbb{R}^n}K(x-\xi,t)\varphi(\xi)\,\mathrm{d}\xi+\int_0^t\mathrm{d}\tau\int_{\mathbb{R}^n}K(x-\xi,t-\tau)f(\xi,\tau)\,\mathrm{d}\xi, \label{eq:poisson-nd}
\end{align}
这里
\begin{align}
K(x,t)=\begin{cases}
\frac{1}{(4\pi a^2 t)^{n/2}}e^{-\frac{|x|^2}{4a^2t}}, & t>0,\\
0, & t\leqslant 0.
\end{cases} \label{eq:poisson-kernel-nd}
\end{align}

\begin{theorem}
设 $\varphi(x)\in C(\mathbb{R}^n)\cap L^\infty(\mathbb{R}^n)$，$f(x,t)\in C^{0,1}(\mathbb{R}_+^{n+1})$，则 Poisson 公式 \eqref{eq:poisson-nd} 表示的函数是初值问题 \eqref{eq:heat-cauchy-nd} 的有界 $C^\infty(\mathbb{R}_+^{n+1})$ 解。
\end{theorem}










\end{document}